%!TeX root = principal.tex
%!TeX encoding = utf-8
As redes de cápsulas foram propostas com a promessa de equivariância: ao representar explicitamente a pose das entidades, seriam mais robustas a transformações ausentes do treino e precisariam de menos dados. Desde então, a augmentação de dados deu às redes convolucionais e de atenção boa parte dessa tolerância. Este trabalho verifica se a vantagem das cápsulas ainda existe. Comparamos o Efficient-CapsNet, a ResNet-18 e o DeiT-Tiny, treinados do zero com hiperparâmetros, partições e sementes idênticos, em CIFAR-10 e em uma versão binária do D-Fire, conjunto de imagens de fogo e fumaça. A grade fatorial soma 162 treinos, com três regimes de augmentação, três frações de dados e três sementes, e cada modelo é avaliado sob rotações grandes e sob uma composição de corrupções não vistas. Os resultados não sustentam as promessas. Sob rotação, a retenção de acurácia do modelo capsular não supera a da ResNet-18 em nenhum dos dois conjuntos, e, no CIFAR-10, a ResNet-18 tem maior acurácia absoluta em todos os regimes com augmentação. A menor queda de acurácia do modelo capsular ao reduzir os dados é de cerca de um ponto percentual e não é significativa. A augmentação amplia a distância em vez de reduzi-la: no CIFAR-10, a desvantagem capsular em acurácia de teste limpa, em média sobre as frações de dados, passa de 4,0 pontos percentuais no regime sem augmentação para 10,4 no regime forte. O Efficient-CapsNet se destaca em outros aspectos: é a arquitetura mais estável entre sementes, sua confiança identifica bem os próprios erros em dados limpos, e no D-Fire fica a poucos pontos da melhor linha de base com um extrator dezenas de vezes menor. O trabalho mostra ainda que a razão de retenção, métrica usual de robustez, favorece modelos fracos, e que composições de muitas corrupções levam todos os modelos ao nível do acaso.